
Este trabalho tem por objetivo o desenvolvimento de um agente inteligente utilizando redes neurais de aprendizado profundo para jogar o pôquer \textit{Texas Hold'Em}. Tendo em vista que o jogo possui uma quantidade quase infinita de estados e acesso incompleto a informações, agentes por aprendizado tendem a cair em ótimos locais e ficarem estagnados sem evoluir. Este trabalho utilizou da técnica de \gls{DQN} e algoritmos de predição \textit{Monte Carlo}, além de modelagem de recompensa e de um motor base 15 capaz de executar desempates de maneira robusta, usufruindo da \gls{API} \textit{TensorFlow.js} usando a ferramenta \textit{Node.js} para desenvolvimento. Os gráficos gerados mostram o sucesso do agente, demonstrando o ganho de fichas de maneira concisa, perdendo menos do que ganhando, exibindo os locais em que o agente explorava e onde o mesmo explotava para receber o máximo possível de fichas. Tornando assim a explotação uma tática viável para agentes que jogam o \textit{Texas Hold'Em}, já que eles conseguiam explorar as falhas dos rivais, também dispensando a necessidade de supercomputadores. A utilização de \textit{JavaScript} permite a implementação do agente por meio de uma arquitetura assíncrona, que permite integração de maneira prática e eficaz.

% Separe as palavras-chave por ponto e vírgula todas em caixa baixa
% resumo de 150 a 200 palavras
\palavraschave{Inteligência Artificial; Texas Hold'Em; Deep Q-Network; Modelagem de Recompensas; TensorFlow.js}